\section{The Ross--Selinger approximation compiler (conditional)}
\label{sec:ross_selinger}

Ross and Selinger give an algorithm that, given an angle $\theta$ and a
precision $\epsilon$, produces a Clifford+T circuit $\epsilon$-approximating
$R_z(\theta)$ with $T$-count $2\log_2(1/\epsilon) + O(\log\log(1/\epsilon))$,
and which is $T$-count optimal relative to a Diophantine factoring oracle.
Formalizing it would supply the logarithmic bound that
Lemma~\ref{lem:lemma12_constants} is missing.

This track is the least complete part of the project, and the shape of what is
and is not proved matters:

\begin{itemize}
  \item The \emph{search architecture} is complete and proved: candidate
    generation, the epsilon region, the Diophantine interface, the completion
    to a unitary, and the reduction of correctness to two named inputs.
  \item Both headline theorems --- soundness
    (Theorem~\ref{thm:rs_sound}) and relative optimality
    (Theorem~\ref{thm:rs_optimal}) --- are stated and their proofs are
    assembled, but they inherit \texttt{sorry} from exactly two open
    mathematical inputs: Lemma~\ref{lem:selinger75} and
    Lemma~\ref{lem:ma_bound}.
  \item Termination is \emph{not} claimed.  Both theorems are conditional in
    the paper's own sense: they say what holds \emph{if} the fuelled search
    returns a circuit.  No factorization, prime-distribution, or termination
    assumption enters.
\end{itemize}

\subsection{Setup}

\begin{definition}
  \label{def:rs_approx}
  \lean{TwoControl.RossSelinger.IsRzApproxCircuit}
  \leanok
  $C$ \emph{$\epsilon$-approximates} $R_z(\theta)$ when
  $\lVert R_z(\theta) - \operatorname{eval}(C)\rVert \le \epsilon$ in operator
  norm.  Note this track uses the operator norm, not the Hilbert--Schmidt
  distance of Section~\ref{sec:circuits}.
\end{definition}

\begin{definition}
  \label{def:candidate}
  \lean{TwoControl.RossSelinger.RSCandidate}
  \leanok
  A \emph{candidate} at level $k$ is a $u \in \mathbb{D}[\omega]$ that lies in
  the Ross--Selinger epsilon region for $(\theta,\epsilon)$, whose
  $\sqrt2$-conjugate lies in the closed unit disk, and whose \emph{least}
  denominator exponent is exactly $k$.  Ordering candidates by least
  denominator exponent, rather than by a chosen presentation, is what makes the
  optimality argument possible.
\end{definition}

\begin{lemma}[Candidates exist]
  \label{lem:grid_exists}
  \lean{TwoControl.RossSelinger.grid_candidate_exists}
  \leanok
  For every $\theta$ and every $\epsilon > 0$ there is $u \in \mathbb{D}[\omega]$
  in the epsilon region for $(\theta,\epsilon)$ whose $\sqrt2$-conjugate lies in
  the closed unit disk.
\end{lemma}

\begin{proof}
  \leanok
  \uses{def:candidate, def:dyadic}
  This is the Ross--Selinger grid lemma: the epsilon region is a rectangle of
  area bounded below in terms of $\epsilon$, and the two-dimensional grid
  $\mathbb{D}[\omega]$ with its conjugate constraint has a point in any such
  region.  The large-$\epsilon$ and small-$\epsilon$ regimes are handled
  separately.
\end{proof}

\begin{lemma}[Executable enumeration]
  \label{lem:grid_bounded}
  \lean{TwoControl.RossSelinger.boundedGridCandidatesAtLevel_sound}
  \leanok
  Scanning an explicit integer coordinate box at a fixed denominator level
  yields only genuine candidates, and every admissible coordinate in the box is
  actually emitted.
\end{lemma}

\begin{proof}
  \leanok
  \uses{def:candidate}
  Soundness is a direct check on each emitted coordinate.  The completeness
  clause is bounded: it is stated relative to the box, so it does not yet give
  full fixed-level completeness.  Closing that gap is a matter of supplying box
  bounds large enough to contain every admissible coordinate.
\end{proof}

\begin{lemma}[Diophantine interface]
  \label{lem:norm_eq}
  \lean{TwoControl.RossSelinger.normEquation_of_solves_completionXi}
  \leanok
  Write $\xi(u) = 1 - \bar u u$.  If $t \in \mathbb{D}[\omega]$ satisfies
  $\bar t t = \xi(u)$, then $(u,t)$ satisfies the norm equation of the unitary
  completion $\operatorname{completionMatrix}(u,t)$.
\end{lemma}

\begin{proof}
  \leanok
  \uses{def:candidate, thm:kmm_completion}
  Direct computation with the norm form.  The role of this lemma is to isolate
  the number-theoretic step --- solving $\bar t t = \xi(u)$ over
  $\mathbb{D}[\omega]$, which is where factorization is needed --- from
  everything else.  Two companion facts complete the interface, both proved: a
  concrete $\mathbb{Z}[\omega]$ solution of the norm equation lifts to a
  solution over $\mathbb{D}[\omega]$, and with exact factorization available the
  solvability predicate is decided, which is how the oracle contract is
  modelled.
\end{proof}

\begin{lemma}[A completed candidate approximates]
  \label{lem:completed_approx}
  \lean{TwoControl.RossSelinger.CompletedCandidate.synthesize_approximates}
  \leanok
  If a candidate for $(\theta,\epsilon)$ is completed by a norm-equation
  solution, the resulting circuit $\epsilon$-approximates $R_z(\theta)$.
\end{lemma}

\begin{proof}
  \leanok
  \uses{lem:norm_eq, def:rs_approx, def:candidate}
  Membership in the epsilon region bounds the operator distance between
  $R_z(\theta)$ and the completion matrix; the norm equation guarantees the
  completion is unitary, and the candidate's circuit evaluates to it.
\end{proof}

\subsection{The two open inputs}

\begin{lemma}[Selinger Lemma 7.5, the $2k-2$ classification]
  \label{lem:selinger75}
  \lean{TwoControl.RossSelinger.selinger_lemma_7_5}
  \leanok
  For a completion matrix with top-left entry $u$ of least denominator exponent
  $k$, one of the two branches $U$ or $TUT^\dagger$ has $T$-count at most
  $\operatorname{rossLevelTCount}(k)$, which is $0$ for $k = 0$ and $2k-2$
  otherwise.
\end{lemma}

\begin{proof}
  \uses{lem:ma_bound, thm:kmm_completion}
  \emph{Proof pending} (\texttt{sorry}).  All the plumbing around it is proved:
  a completion has the same least $\omega$-denominator exponent as its entry,
  the $TUT^\dagger$ branch is converted into the adjusted completion
  $\omega t$ with its norm equation, and either branch is packaged for the
  search layer.  What remains is the finite Ma/KMM residue classification that
  selects the branch, which is meant to come from
  Lemma~\ref{lem:ma_bound}.
\end{proof}

\begin{lemma}[Matsumoto--Amano normal form, Figure 2 witness]
  \label{lem:ma_bound}
  \lean{TwoControl.RossSelinger.ma_normal_form_figure2_synthesis_bound}
  \leanok
  Let $U$ be a one-qubit unitary with entries in $\mathbb{D}[\omega]$, least
  $\omega$-denominator exponent $k$, and residue $R$ at level $k$ matching a
  node of the Figure 2 automaton valid at level $k$.  Then $U$ is exactly a
  Clifford+T circuit of $T$-count at most $2k - \operatorname{tOffset}(\text{node})$.
\end{lemma}

\begin{proof}
  \uses{thm:kmm_exact, def:sde}
  \emph{Proof pending} (\texttt{sorry}).  The Matsumoto--Amano syllable
  semantics, the Figure 2 automaton, the residue bookkeeping, and the extraction
  of the synthesis statement from a witness are all formalized; the pending step
  is the construction of the witness itself, i.e.\ that the automaton path
  selected by the residue really does reduce $U$ one syllable at a time while
  decreasing the denominator exponent.  Note the syllable words act as
  \emph{reduction prefixes}: from a circuit for $N^\dagger U$ one obtains a
  circuit for $U$ by prefixing $N$.
\end{proof}

\subsection{The conditional theorems}

\begin{theorem}[Conditional soundness]
  \label{thm:rs_sound}
  \lean{TwoControl.RossSelinger.rossSelingerSearch_sound_if_returns}
  \leanok
  If the fuelled Ross--Selinger search returns a circuit $C$ on input
  $(\theta,\epsilon)$, then $C$ $\epsilon$-approximates $R_z(\theta)$.
\end{theorem}

\begin{proof}
  \uses{lem:completed_approx, lem:grid_bounded, lem:norm_eq, lem:selinger75}
  \emph{Proof pending} (\texttt{sorry}), inherited from
  Lemma~\ref{lem:selinger75}.  The argument itself is complete: a returned
  circuit comes from a completed candidate, candidate batches carry the
  epsilon-region certificate by Lemma~\ref{lem:grid_bounded}, the solver law
  turns the returned $t$ into a norm-equation certificate, and
  Lemma~\ref{lem:completed_approx} concludes.  Nothing about when the search
  returns is claimed.
\end{proof}

\begin{definition}
  \label{def:rs_optimal}
  \lean{TwoControl.RossSelinger.IsOptimalRzApproxByTCount}
  \leanok
  $C$ is \emph{$T$-count optimal} for $(\theta,\epsilon)$ when it
  $\epsilon$-approximates $R_z(\theta)$ and no circuit with that property has a
  smaller $T$-count.
\end{definition}

\begin{theorem}[Conditional relative optimality]
  \label{thm:rs_optimal}
  \lean{TwoControl.RossSelinger.rossSelingerOracle_optimal_if_returns}
  \leanok
  Given a factoring oracle and the competitor bridge, if the oracle variant of
  the search returns a circuit $C$, then $C$ is $T$-count optimal for
  $(\theta,\epsilon)$.
\end{theorem}

\begin{proof}
  \uses{thm:rs_sound, def:rs_optimal, lem:selinger75, lem:ma_bound}
  \emph{Proof pending} (\texttt{sorry}), inherited from
  Lemma~\ref{lem:selinger75}.  The competitor side is factored along the
  paper's argument, with the paper-side inputs collected into an explicit
  hypothesis bundle: any competing circuit is put into completion-matrix normal
  form, the resulting scaled-grid candidate is exposed by the enumerator at its
  least denominator exponent, and the denominator-exponent theorem gives that
  competitor's $T$-count lower bound.  The search-order argument --- that once a
  candidate returns at level $k$, the oracle cannot have skipped a completable
  candidate at a lower level --- is proved.  Combining it with the returned
  branch bound of Lemma~\ref{lem:selinger75} gives optimality.
\end{proof}
